% Chapter 36: Polarization and Light in Matter
\chapter{Polarization and Light in Matter}

\step{A Counterintuitive Opening}

Imagine a water purifier. Water flows through one filter and then another. Insert a denser filter between them and the output can only decrease, never increase. More filtering blocks more: everyone's intuition about filters.

Now replace water with light. Find polarized sunglasses, as most sunglasses on the market are, and an old LCD monitor or LCD phone. The screen emits light that is already ``organized,'' as this chapter will explain. Rotate a sunglass lens slowly in front of it. At one angle the screen looks normal; ninety degrees later it is nearly black. The lens has filtered out its light.

Now try the impossible. Keep that first lens at the black-screen angle, put another polarizer in front, and rotate the second. With two lenses, the screen stays black however you turn it. Sensible: two layers of filtering.

Then comes the miracle: insert the second lens \textbf{diagonally between the other two}, with the three orientations successively vertical, forty-five degrees, and horizontal. Light emerges through the previously black screen! Add another light-blocking filter and light gets through.

No magic equipment is needed; any two polarizers can repeat it. This slaps filtering intuition in the face: polarizers plainly do more than ``block a fraction.'' We must discover the unnoticed extra ``knob'' light carries and how it explains blue skies, rainbows, and liquid-crystal screens.

A fair word for intuition: finding this surprising shows your common sense is healthy. It comes from macroscopic water and air filters, which encounter only one relevant dimension, how much material to stop. Light is not a stream of impurities awaiting filtration. Experiment will expose a hidden dimension that ordinary filters never had a chance to show you. Your intuition was not wrong; it had never worked this particular shift.

\step{Make a Guess First}

Write down your guess before continuing. Which account sounds closest to the truth?

First: polarizers are directional color filters. Stacked dyes darken the light, and the third layer happens to adjust the color back.

Second: light is a stream of little particles flying straight, and polarizers are fences whose gaps the particles pass through. Three fences happen to align into a passage.

Third: besides direction, color, and brightness, light has an extra property that a polarizer reads and writes. The third layer does not simply ``let light through''; it first \textbf{rewrites} this property so the next layer can transmit it.

Record your choice for checking at the end. Remember the water-filter prediction, ``another blocking layer means less transmission.'' If experiment breaks it, the casualty is not merely an account of polarizers, but the deeper image of light fully described as something traveling along a trajectory.

One hint before placing bets, without an answer: the first account has a fatal flaw. No stack of color filters brightens two completely black layers by adding another; dyes subtract light, not create it. The second is vividly imaginable, so remember it particularly: some of the finest minds centuries ago placed that bet, whose verdict arrives midway through the chapter. The third remains empty words for now. What extra property, and what does rewriting mean? Empty words must pass through experiment to become physics. Let us experiment.

\step{Hands-on Experiment}

\begin{experiment}[Three Filters and the Screen's Dark Squares]
Materials: a phone with an LCD, though OLED works less well, and polarized sunglasses. For a second polarizer, try 3D cinema glasses, another pair of polarized sunglasses, or even film removed from an old calculator.

Step one: display a plain white phone screen, such as an empty notes page. Hold the polarizing lens against it and rotate through a full turn. Brightness reaches two maxima and two minima. Stop at the darkest angle, where the lens's ``direction'' is perpendicular to the screen's.

Step two: insert a second polarizer diagonally between lens and screen, roughly forty-five degrees to both. The black screen brightens to about one-eighth its unobstructed brightness. Remove the middle sheet and darkness returns; replace it and light returns. Repeat until the surprise settles into a question.

% BEGIN TRANSLATOR SAFETY NOTE
\textbf{Translator safety note:} Observe sky away from the Sun, not the Sun itself. Ordinary sunglasses and polarizing films are not safe solar-viewing filters, however dark they appear. See \href{https://science.nasa.gov/eclipses/safety/}{NASA solar-viewing guidance}.
% END TRANSLATOR SAFETY NOTE

Step three: take the lens outside. Look at sky ninety degrees from the Sun and rotate it; that patch noticeably darkens and brightens, showing that skylight is also organized. Look down at glare from water or asphalt and rotate to an angle where the reflection nearly disappears, revealing things beneath the water and road surface. At strongest suppression, the angle involving Sun, water, and your eyes is fifty-something degrees. We will name it later.
\end{experiment}

Record three observations: two brightness cycles per full rotation, not one; insertion of a middle layer brightens blackness; and water glare almost vanishes at a particular angle. Each is a clue.

Warm up with a little reasoning. Why two maxima and minima per turn? If the lens recognizes direction, a half-turn restores its original state. Rotating a glass sheet by one hundred eighty degrees does not turn it upside down; it has no top or bottom. Brightness must therefore repeat after half a turn. This small symmetry argument suggests an axis-like property rather than an arrow-like one: an axis repeats after half a turn, an arrow after a full turn. Keep the distinction; it will recur.

\step{The Old Model Runs into Trouble}

Inventory the old models and see where each sticks.

Chapter 33's ray is fully specified by position and direction. It explains shadows, reflection, and refraction, but its record contains no entry for ``rotated ninety degrees about its axis.'' A line flying straight remains the same line after such a rotation. Rays therefore predict that rotating a polarizer changes nothing. Experiment says otherwise.

Chapter 35's waves improve matters, but not enough, because there are two kinds. In longitudinal waves such as sound, vibration parallels propagation; in transverse waves such as waves on a rope, it is perpendicular. Longitudinal light would also have no orientation around its axis: rotate a compressing spring and its wave remains the same. Airborne sound sounds the same whether the speaker lies horizontally or vertically. Thus longitudinal light cannot explain the rotating lens either.

Only one route remains: light is transverse, and transverse waves can vibrate in different directions. A horizontally traveling transverse wave can oscillate up-down, side-to-side, or diagonally. Rotation about propagation really changes its state. But trouble immediately returns: natural sunlight or lamplight through one polarizer always has the same brightness, exactly half, whatever the lens angle. If natural light had a definite vibration direction, why no preferred angle? Stranger still, a perpendicular second polarizer extinguishes the first one's output, while inserting a diagonal layer restores it. Filtering models cannot cope: filters delete, not rewrite.

\begin{mainline}
The difficulty is this: \textbf{besides direction, color, and brightness, light carries an orientation about its axis. Natural light is a jumble of orientations, and a polarizer is not a sieve but a machine that projects and rewrites this orientation.} We need a name for it.
\end{mainline}

\step{Inventing New Concepts}

Build the concept yourself, starting mechanically. Tie a long rope to a doorknob and shake the other end up and down, launching an up-down transverse wave. Shake side to side for another transverse pattern. These are independent vibration modes on the same rope. Thread it through a vertical fence gap: up-down motion passes because the rope can move vertically within the gap, while sideways motion is blocked. For rope waves, the fence truly is a sieve.

Light goes deeper. Electromagnetism says it consists of electric and magnetic fields varying in space and time, jointly satisfying Maxwell's equations and propagating as a wave. The electric field's oscillation direction is its \textbf{polarization direction}. It resembles a rope's vibration mode because both are transverse. More importantly, it differs because the rope is tangible matter, while light's electric field oscillates without any material medium.

Before classifying polarization, establish a basic skill: any linear vibration can be decomposed into vertical and horizontal parts. This is vector decomposition, not a new optical law. A rope vibrating at forty-five degrees is equivalent to equal, in-phase vertical and horizontal vibrations occurring together. Conversely, those two parts combine into a diagonal vibration. This decomposition and recombination is the whole secret of ``rewriting'': the polarizer splits the incident vibration into components parallel and perpendicular to its transmission axis, blocks the perpendicular part, and passes the parallel part. Emerging light therefore takes the transmission axis's direction. Now classify the possibilities:

\textbf{Linear polarization}: the electric field oscillates along a fixed line. Vertical and horizontal rope waves are relatives. An ideal polarizer produces linear polarization aligned with its transmission axis.

\textbf{Natural, or unpolarized, light}: sunlight and lamplight arise from independent radiation by countless atoms. Each short segment has a random polarization direction, with no direction predominating in the mixture. It is rotationally symmetric, so every angle of the first polarizer gives the same average, one half. This explains ``one layer halves it.''

\textbf{Circular polarization}: imagine a constant-magnitude electric field rotating uniformly like a propeller. As light advances, the field arrow traces a spring-like shape along the route. This sounds exotic but decomposes simply: take equal vertical and horizontal vibrations offset by a quarter-cycle, one at its maximum when the other crosses zero. Together they rotate. Cinema 3D glasses exploit this.

\textbf{Elliptical polarization}: the more general version of circular polarization, with field magnitude changing as its direction rotates so its tip traces an ellipse. Linear and circular cases are special limits: flatten the ellipse to a line or round it into a circle.

There is also intermediate \textbf{partially polarized light}, such as skylight, mostly polarized with a smaller disordered part. A lens darkens it but cannot make it completely black.

Circular polarization is not merely mathematical curiosity; it already sits on your nose. Modern 3D cinemas use two projectors, one with left-handed circular polarization and one right-handed. Each eyeglass lens passes only its assigned handedness, sending different images to the two eyes for the brain to assemble into depth. The practical advantage over linear polarization is head tilt: ``vertical'' and ``horizontal'' tilt with your head, mixing the images, while clockwise and counterclockwise remain distinct. Old red-blue 3D glasses selected color; this generation selects rotational handedness, upgrading which property carries the separation.

Nature used polarization before engineers. Bees' compound eyes detect skylight polarization, whose clear-sky pattern is uniquely determined by the Sun's position. It is strongest ninety degrees from the Sun, forming a solar-centered direction map. Bees follow it home even under clouds or when mountains hide the Sun. The darkening and brightening sky in step three is an edge of their map.

Polarization reveals half the solution to the three-layer miracle: the middle sheet does not merely filter again; it \textbf{rewrites} vertical light into forty-five-degree light. Diagonal vibration includes a horizontal component, giving the last sheet something to pass. The remaining half, how much rewriting occurs, belongs to the mathematical translator.

One last concept concerns the title's second half, light in matter. Polarization matters partly because it participates directly in every encounter between light and material. Matter contains charged particles, negative electrons and positive nuclei. A passing optical electric field drives them to oscillate, and oscillating charges themselves radiate electromagnetic waves. Everything light does in matter, refraction, dispersion, absorption, scattering, and birefringence, is a variation on ``light drives charges; charges reradiate'':

\textbf{Origin of refractive index}: in transparent glass or water, incident and reradiated waves combine to delay crest positions, effectively slowing propagation. Chapter 33's index \(n\) microscopically measures this delay.

\textbf{Dispersion}: electrons are not completely free but tethered to nuclei like spring-mounted balls with preferred natural frequencies. The closer the incident frequency to those frequencies, the stronger the response and delay. Violet lies closer than red to glass electrons' natural frequencies, so it is delayed more, has higher index, and bends farther. A prism separates white light because glass delays different colors differently.

\textbf{Absorption}: when light matches a charge's natural frequency, like a well-timed swing push, it drives strong oscillation. Collisions transfer this energy to neighboring particles, genuinely converting optical energy to internal energy. That is absorption. Glass's relevant natural frequencies lie in the ultraviolet, so visible light cannot drive them and glass is transparent.

\textbf{Scattering}: air molecules also capture light energy and radiate it in every direction. For molecules much smaller than wavelength, scattering efficiency rises steeply with frequency, spreading blue-violet much more strongly than red. The blue sky is light scattered in all directions, with blue especially abundant. The setting Sun looks red because much of its blue was scattered away en route, leaving the remainder to reach your eyes directly.

\textbf{Birefringence}: the preceding account assumed identical properties in all directions. Crystals such as calcite have different internal ``spring stiffnesses'' along different directions. One polarization encounters less resistance, its perpendicular counterpart more. The same crystal therefore has two refractive indices for different polarizations and splits one beam into two. Put calcite over a printed character and see two displaced copies. Polarization becomes a probe rather than merely something observed, detecting matter's internal directionality.

Metals take driving and reradiation to an extreme. Their electrons are nearly untethered, free to respond immediately to light of any frequency, and they respond extraordinarily well. Their collective reradiation exactly cancels the incident wave inside while sending a strong reflected wave back into the air. Thus metals are opaque and shiny. One charge-relay law gives transparency and refraction for bound electrons, reflection and metallic luster for free ones.

\step{The Model in One Sentence}

\begin{oneline}
Light is transverse, and its electric-field oscillation direction is polarization. A polarizer does not delete light; it projects polarization onto its transmission axis. Everything light does inside matter, slowing, separating colors, being absorbed, or scattering into the sky, arises from how strongly and how promptly driven charges reradiate.
\end{oneline}

Read this twice in parts. First focus on projection: a polarizer takes the electric vector's component along its transmission axis rather than simply subtracting, which lets three sheets revive light. Then focus on driving and reradiation: light does not burrow through matter searching for a path, but participates in a charge relay whose pace depends jointly on frequency and material structure.

\step{The Mathematical Translator}

Intuition is ready; translate it into equations.

For one polarizer, let the incident linearly polarized field amplitude be \(E_0\), with angle \(\theta\) to the transmission axis. Elementary vector geometry gives the transmitted component \(E_0\cos\theta\); the perpendicular component is absorbed. Intensity is proportional to field amplitude squared, akin to the quadratic energy dependence of a more vigorously shaken rope. Thus transmitted intensity is
\[
  I = I_0 \cos^2\theta .
\]
This is \textbf{Malus's law}. The left answers ``how bright is the output?'' On the right, \(I_0\) is the incoming capital and \(\cos^2\theta\) results from projection followed by squaring amplitude into intensity.

\begin{center}
\begin{tikzpicture}[scale=1.7,>=Stealth]
  \draw[->] (0,0) -- (2.8,0) node[right] {Transmission axis};
  \draw[->] (0,0) -- (0,1.8);
  \draw[->,very thick] (0,0) -- (40:2.3) node[above right=-2pt] {$E_0$};
  \draw[->,very thick,blue!60!black] (0,0) -- ({2.3*cos(40)},0) node[below] {$E_0\cos\theta$};
  \draw[dashed] (40:2.3) -- ({2.3*cos(40)},0);
  \draw (0.8,0) arc[start angle=0,end angle=40,radius=0.8];
  \node at (20:1.1) {$\theta$};
\end{tikzpicture}
\end{center}

Check special cases immediately. At \(\theta=0\), \(\cos^2\theta=1\): all passes, since directions already align. At \(\theta=90^\circ\), \(\cos^2\theta=0\): all vanishes, since vertical vibration has zero horizontal component, explaining the first experiment's blackness. At \(\theta=45^\circ\), \(\cos^2\theta=1/2\): half passes. For unpolarized light, averaging \(\cos^2\theta\) over equally represented directions gives exactly \(1/2\), so any first sheet halves it, matching experiment.

Now calculate the miracle. Unpolarized intensity \(I_0\) passes the first sheet as \(I_0/2\), vertically polarized. The middle sheet at \(45^\circ\) gives \((I_0/2)\times\cos^2 45^\circ = I_0/4\), diagonally polarized. The horizontal third sheet differs by another \(45^\circ\), giving \((I_0/4)\times\cos^2 45^\circ = I_0/8\). Remove the middle and vertical meets horizontal: \(I_0/2\times\cos^2 90^\circ=0\). One-eighth revival matches the experiment's approximate brightness. The purifier intuition is officially defeated: cascaded polarizers are a rotational relay, not merely multiplying filters.

Estimate with real numbers what \(1/\lambda^4\) does at sunset. Earth's radius is about \(6400\,\mathrm{km}\), while effective atmospheric thickness is only about \(8\,\mathrm{km}\). Noon sunlight crosses those \(8\,\mathrm{km}\) almost vertically; sunset light skims in along the ground, traveling tens of times farther through the same air, over thirty times by rough geometry. Every length of travel pays a wavelength-dependent scattering tax: blue's rate is over four times red's. Tens of times the distance means repeated collection, nearly exhausting blue while leaving substantial red. The noon Sun is dazzling white, the sunset Sun softly red. Same Sun, same air; geometry enlarges only the path length by an order of magnitude.

\begin{mathbridge}
The angle of suppressed water glare is calculable too. At an interface, polarization parallel to the incidence plane reflects differently from polarization perpendicular to it. At one special angle, the parallel reflection coefficient vanishes. Reflected and refracted rays are then perpendicular, and the reradiating charges oscillate directly along the reflected direction. A transverse wave cannot propagate along its own vibration direction, so that reflection disappears. Refraction and this right-angle condition give
\[
  \tan\theta_{\mathrm B} = \frac{n_2}{n_1}.
\]
This is the \textbf{Brewster angle}. Water has \(n_2=1.33\) and air \(n_1=1\), giving \(\theta_{\mathrm B}=\arctan 1.33\approx 53^\circ\), the experiment's fifty-something degrees. Glass with \(n_2=1.5\) gives about \(56^\circ\).

\begin{center}
\begin{tikzpicture}[scale=1.1,>=Stealth]
  \draw[thick] (-4,0) -- (4,0) node[right] {Water surface};
  \fill[blue!8] (-4,-1.6) rectangle (4,0);
  \draw[dashed] (0,0) -- (0,2.6);
  \draw[->,very thick] ({-2.6*tan(53)},2.6) -- (0,0) node[midway,left,xshift=-2pt] {Incident};
  \draw[->,very thick] (0,0) -- ({2.6*tan(53)},2.6) node[midway,right,xshift=2pt] {Reflected (linear)};
  \draw[->,very thick,blue!50!black] (0,0) -- ({1.9*tan(37)},-1.9) node[midway,right,xshift=2pt] {Refracted};
  \draw (0,1.1) arc[start angle=90,end angle=143,radius=1.1];
  \node at (118:1.55) {$\theta_{\mathrm B}$};
  \draw (37:0.7) arc[start angle=37,end angle=-53,radius=0.7];
  \node at (-10:1.05) {$90^\circ$};
  \node[align=center] at (0,-2.3) {At Brewster incidence, reflected and refracted rays are perpendicular};
\end{tikzpicture}
\end{center}

Check an extreme: \(n_2=n_1\) gives \(\theta_{\mathrm B}=45^\circ\), but no interface then exists, so ``fully polarized reflection'' has no meaning. However elegant the angle, an interface must exist for it to matter.
\end{mathbridge}

\begin{mathbridge}
Dispersion and scattering each have a calculable statement. Glass electrons' natural frequencies lie in the ultraviolet. At lower visible frequency \(\omega\), index rises gently and monotonically with frequency, so violet's index exceeds red's, fixing rainbow order. For particles far smaller than wavelength, scattered intensity grows as frequency to the fourth power, equivalently inverse wavelength to the fourth:
\[
  I_{\text{scat}} \propto \frac{1}{\lambda^4}.
\]
For blue \(\lambda\approx 450\,\mathrm{nm}\) and red \(\lambda\approx 650\,\mathrm{nm}\), the ratio is \((650/450)^4\approx 4.3\): blue scatters over four times more efficiently. This starts the order-of-magnitude account of why the sky is blue rather than violet. Solar spectrum and human sensitivity must also be included before violet finally loses to blue.
\end{mathbridge}

\begin{challenge}
Linear, circular, and elliptical polarization share a notation: write electric-field components in two perpendicular directions as two numbers with their relative phase, forming a two-component Jones vector. Horizontal is \((1,0)\), vertical \((0,1)\), forty-five-degree linear \((1,1)/\sqrt{2}\), and circular \((1,\pm\mathrm{i})/\sqrt{2}\). The imaginary unit \(\mathrm{i}\) records the quarter-cycle offset. Every polarizer or birefringent plate becomes a \(2\times2\) matrix, and an optical train becomes matrix multiplication. Three polarizers' \(I_0/8\) follows in three lines. This notation returns almost unchanged in Chapter 46: polarization is quantum mechanics' simplest two-state system, where \((1,0)\) and \((0,1)\) are renamed basis-state vectors.
\end{challenge}

\step{The Model's Limits}

Every model has a jurisdiction; crossing it produces mistakes. Here are this chapter's boundaries.

Malus's law assumes completely linearly polarized incident light and an ideal polarizer. Real sheets absorb some parallel component too, and natural light is isotropic only statistically. Keep those small corrections in mind when calculating \(I_0/8\).

Retire the fence analogy explicitly. \textbf{Where it fits}: for mechanical transverse waves on ropes, a gap really passes only vibration parallel to it. \textbf{Where it fails}: a polarizer has no gaps. Its axis arises from ordered molecules selectively absorbing field components. More fatally, a fence only deletes and cannot rewrite, predicting that three fences darken further while three polarizers brighten. Thinking only in terms of sieving will mislead you.

Rayleigh scattering's \(1/\lambda^4\) applies only to particles much smaller than wavelength. It is tempting to extrapolate strong blue scattering to clouds: why are their little droplets not blue? Droplets are a dozen or more micrometers across, tens of wavelengths, invalidating the \(1/\lambda^4\) premise. Large particles scatter colors roughly equally and clouds look white. Changing particle scale reverses the conclusion within the same general phenomenon of scattering, a useful counterexample to intuition.

The dispersion rule ``closer to resonance means higher index'' also has limits. At the natural frequency itself, delay is no longer mild and absorption dominates; glass becomes opaque. In anomalous-dispersion regions, frequency dependence disrupts the simple rainbow intuition, alongside strong absorption bands.

Birefringence requires directional material properties. Ordinary glass is isotropic with one index. But bend a plastic ruler strongly and stress temporarily creates directionality. Between crossed polarizers it reveals colored stress fringes, the method engineers use to photograph stress in transparent models.

% BEGIN TRANSLATOR SAFETY NOTE
\textbf{Translator safety note:} Polarization or dark tint does not make a lens safe for looking directly at the Sun. Do not use these films or sunglasses for solar observation. See \href{https://science.nasa.gov/eclipses/safety/}{NASA solar-viewing safety guidance}.
% END TRANSLATOR SAFETY NOTE

A shopping misconception deserves separate mention: polarizers are not synonymous with sunglasses. Ordinary dark lenses dim every component without selecting polarization; only polarized lenses contain the projection machine. You already know the test: rotate a lens while viewing an LCD. Strong brightness changes indicate polarization; no change indicates merely dark glass. Both block sunlight, but only polarizers suppress water and road glare, because that light is not merely too bright but too organized.

Finally, do not picture polarization as a tiny arrow attached to a flying photon. It is a collective property of electromagnetic-field oscillation. What a single photon's polarization means requires probability and belongs to the quantum section.

\step{Explaining the Opening Phenomenon}

Now settle the opening's four accounts.

Three sheets revive light: the first projects natural light vertically, \(I_0/2\); the middle projects it diagonally, \(I_0/4\); the last horizontally, \(I_0/8\). Each rewrites polarization rather than merely deleting, so another sheet can restore transmission. The third proposed explanation was right: polarization is an independent optical property, and polarizers read and write it.

The screen darkens under a rotated lens: each LCD pixel uses a backlight that first passes a polarizer, then a liquid-crystal layer. Voltage-controlled molecular alignment rotates polarization like the middle layer of the experiment. A final polarizer determines whether the pixel is bright or dark. Your screen is essentially hundreds of thousands of miniature three-polarizer arrangements, replacing the middle sheet with electrically controlled liquid crystal. Hence its naturally polarized output, the foundation of the LCD industry.

Water glare disappears: sunlight incident at Brewster's angle, about \(53^\circ\), reflects almost entirely with horizontal linear polarization. Sunglasses have vertical transmission axes that project it to zero while admitting more than half of the disordered light from underwater. Anglers seeing fish and photographers suppressing shop-window reflections solve the same problem.

Blue skies and red sunsets: atmospheric molecules scatter blue-violet light across the sky about four times more efficiently than red. Your eyes receive this scattered light, making the sky blue. Evening paths are tens of times longer, scattering nearly all blue away before the direct Sun reaches you in orange-red. One mechanism gives blue in one direction and red in another: the reversed-direction check passes.

Settle two incidental accounts too. In 3D cinema, two projectors emit opposite circular handedness, each lens passes one, and head tilt does not mix images because handedness stays distinct. Calcite gives two indices to two polarization directions, making two refracted paths and two copies of one character. The sky patch darkened by a rotating lens contains blue light scattered toward you through ninety degrees. Transverse-wave rules leave almost one polarization direction, which the lens suppresses. Bees read this map daily; now you can too.

\step{Thought Experiments and Challenges}

\begin{thought}[One Photon at a Time]
Dim the source until it emits on average one photon per short interval, already achievable with modern laboratory single-photon sources. Send them individually toward two polarizers separated by angle \(\theta\). Each photon either passes whole or disappears whole, never transmitting one-eighth of a photon. Malus's law becomes probabilistic: each photon's transmission probability is \(\cos^2\theta\), and many photons accumulate intensity \(I_0\cos^2\theta\).

Reconsider three sheets. Inserting the middle one at \(45^\circ\) repeatedly redraws a photon's fate. Polarization now resembles not a classical arrow but a record specifying probabilities at the next encounter. This is a smooth bridge to Chapters 43 and 45: quantum states, probability amplitudes, and measurement can all be rehearsed with three inexpensive plastic sheets. Remember the lesson: the polarizer does not coarsely disturb the photon, and probability is not poor workmanship but nature's own accounting.
\end{thought}

\begin{puzzles}
\item Rotate two ideal polarizers from \(0^\circ\) to \(90^\circ\); transmitted intensity falls from \(I_0\) to \(0\). Now use three, with first and third fixed perpendicular and the middle rotating from \(0^\circ\) to \(90^\circ\). Does intensity change monotonically or rise and then fall? At which middle angle is it brightest, and by how much? (Hint: write both \(\cos^2\) factors; the two angles always sum to \(90^\circ\). Calculate at \(0^\circ\), \(45^\circ\), and \(90^\circ\) before guessing.)
\item On a rainy night, wet asphalt reflects headlights with fierce glare. A passenger suggests polarized sunglasses; the driver objects that the road would become hard to see. Who has the better argument? (Hint: what is the glare's polarization, and what about diffuse road light carrying actual surface information? Think further: what information would disappear if every driver wore them?)

% BEGIN TRANSLATOR SAFETY NOTE
\textbf{Translator safety note:} Do not wear tinted sunglasses at night if they restrict your vision; glare reduction does not justify losing road visibility. See \href{https://www.gov.uk/guidance/the-highway-code/rules-for-drivers-and-motorcyclists-89-to-102}{the Highway Code, Rule 94}.
% END TRANSLATOR SAFETY NOTE
\item Outside Earth's atmosphere, astronauts see black rather than blue sky; the Moon's sky is also black. Explain the contrast with Earth and predict how Mars's thin, dusty atmosphere changes its sky color. (Hint: blue sky requires scatterers, and those scatterers must be smaller than wavelength. Check which condition is missing on each world.)
\end{puzzles}
